\SourcePage{79}
\SGAsection{3}{凝聚层 $\underline D_X^i$}
通常不能\Term{prorepresentable}{可预表}、却仍容许\Term{formal-versal-object}{形式万有对象}的良好\Term{functor}{函子}之一，是\Term{deformation-functor}{形变函子}。为此，我们先回顾关于\Term{commutative-ring}{交换环}\Term{group-extension}{扩张}及其\Term{obstruction-theory}{障碍理论}的若干基本事实。本节没有新的内容（或许通过\Term{site}{景}上的 Grothendieck \Term{cohomology}{上同调}来作系统处理这一点除外）。本节的一个目的，是使本讲尽可能自足。（事实上，这里证明的内容比后面所需的更多。）基本事实陈述于\Term{theorem}{定理} 3.12。全文中，环均指有\Term{identity-element}{单位元}的\Term{commutative-ring}{交换环}。

\SourcePage{80}
给定环 $R$，以 $\mathcal C_R$ 表示 $R$-\Term{algebra}{代数}的\Term{category}{范畴}。对 $\mathcal C_R$ 的每个\Term{object}{对象} $A$，令
\[
\operatorname{Cov}(A)=\{\mathcal C_R\text{ 中以 }A\text{ 为终点的所有满射箭头 }\varphi:A'\to A\}.
\]
这显然定义一个\Term{pretopology}{预拓扑}，从而得到一个\Term{site}{景} $\mathcal C_R$。对 $\mathcal C_R$ 的每个\Term{object}{对象} $A$，以 $\mathcal C_{A\mid R}$ 表示 $\mathcal C_R$ 中以 $A$ 为终点的\Term{arrow}{箭头}\Term{category}{范畴}。我们可以且将赋予它由遗忘\Term{functor}{函子} $j_A:\mathcal C_{A\mid R}\to\mathcal C_R$ 诱导的\Term{topology}{拓扑}。注意，$\mathcal C_{A\mid R}$ 有\Term{final-object}{终对象} $(1_A:A\to A)$ 和\Term{initial-object}{始对象} $(i_A:R\to A)$，并且有\Term{fiber-product}{纤维积}和\Term{coproduct}{余积}。今后，在不致混淆时，我们滥用记号，将 $\mathcal C_{A\mid R}$ 的\Term{object}{对象} $(B\to A)$ 简记为 $B$。给定一个\Term{abelian-category}{阿贝尔范畴} $\mathcal C$，$\mathcal C_{A\mid R}$ 上的 $\mathcal C$ 值\Term{sheaf}{层}，按定义是一个\Term{functor}{函子}
\[F:\mathcal C_{A\mid R}^{\circ}\longrightarrow\mathcal C,\]
使得对 $\mathcal C_{A\mid R}$ 中任意\Term{surjectivity}{满射}\Term{arrow}{箭头} $B'\to B$，序列
\[
0\longrightarrow F(B)\longrightarrow F(B')\rightrightarrows F(B'\mathbin{\times_B}B')
\]
均\Term{exact}{正合}。所谓\Term{additive-presheaf}{加性预层}，是指保持\Term{coproduct}{余积}的\Term{presheaf}{预层}，即
\[F(B_1\mathbin{\amalg_R}B_2)=F(B_1)\oplus F(B_2).\]
例如，一个 $A$-模 $M$ 定义\Term{functor}{函子}
\[\operatorname{Der}_R(-,M):\mathcal C_{A\mid R}^{\circ}\longrightarrow(A\text{-mod}),\]
将 $B$ 对应于 $\operatorname{Der}_R(B,M)$，即 $B$ 到 $M$ 的 $R$-\Term{linear}{线性}\Term{derivation}{导子}所成的模；这里通过结构映射 $B\to A$ 将 $M$ 看作 $B$-模。显然，这是 $\mathcal C_{A\mid R}^{\circ}$ 上的\Term{additive-sheaf}{加性层}；$(A\text{-mod})$ 表示 $A$-模\Term{category}{范畴}。$\mathcal C_{A\mid R}$ 上取值于 $(A\text{-mod})$ 的\Term{presheaf}{预层} $F$ 称为“\Term{finite-type}{有限型}”，若每当 $B$ 是\Term{finite-type}{有限型} $R$-\Term{algebra}{代数}时，$F(B)$ 都是\Term{finite-type}{有限型} $A$-模。例如，若 $A$ \Term{noetherian}{诺特}且 $M$ 是\Term{finite-type}{有限型} $A$-模，则 $\operatorname{Der}_R(-,M)$ 是\Term{site}{景} $\mathcal C_{A\mid R}$ 上的\Term{finite-type}{有限型}\Term{additive-sheaf}{加性层}。

再说明一些\Term{homological-algebra}{同调代数}的记号和基本事实。

\textbf{3.1(a)} 本节以下部分，用 $\widetilde{\mathcal C}_{A\mid R}$（相应地，$\mathcal C_{A\mid R}^{\vee}$）表示取值于 $(A\text{-mod})$ 的\Term{sheaf}{层}（相应地，\Term{presheaf}{预层}）\Term{category}{范畴}。有（\Term{left-exact}{左正合}的）\Term{embedding}{嵌入}\Term{functor}{函子}
\[i_{A\mid R}:\widetilde{\mathcal C}_{A\mid R}\longrightarrow\mathcal C_{A\mid R}^{\vee},\]
\SourcePage{81}
以及与 $i_{A\mid R}$ \Term{adjunction}{伴随}的（\Term{exact}{正合}的）\Term{associated-sheaf}{伴随层}\Term{functor}{函子}
\[G_{A\mid R}:\mathcal C_{A\mid R}^{\vee}\longrightarrow\widetilde{\mathcal C}_{A\mid R}.\]
\Term{object}{对象} $B\in\operatorname{ob}(\mathcal C_{A\mid R})$ 诱导\Term{continuous-functor}{连续函子} $j:\mathcal C_{B\mid R}\to\mathcal C_{A\mid R}$，从而给出\Term{commutative-diagram}{交换图}
\[
\begin{tikzcd}
\widetilde{\mathcal C}_{A\mid R}\arrow[r,"\widetilde j"]\arrow[d,"i_{A\mid R}"']&\widetilde{\mathcal C}_{B\mid R}\arrow[d]\\
\mathcal C_{A\mid R}^{\vee}\arrow[r,"j^{\vee}"']&\mathcal C_{B\mid R}^{\vee}
\end{tikzcd}\tag{*}
\]
其中\Term{restriction}{限制}\Term{functor}{函子} $\widetilde j$（以及 $j^{\vee}$）是\Term{exact}{正合}\Term{functor}{函子}，因为\Term{site}{景} $\mathcal C_{B\mid R}$ 就是 $\mathcal C_{A\mid R}$ 中以 $B$ 为终点的\Term{arrow}{箭头}\Term{category}{范畴}，赋予诱导\Term{topology}{拓扑}。\Term{ring-homomorphism}{环同态} $S\to R$ 诱导\Term{continuous-functor}{连续函子} $\beta:\mathcal C_{A\mid R}\to\mathcal C_{A\mid S}$，从而给出\Term{commutative-diagram}{交换图}
\[
\begin{tikzcd}
\widetilde{\mathcal C}_{A\mid S}\arrow[r,"\widetilde\beta"]\arrow[d,"i_{A\mid S}"']&\widetilde{\mathcal C}_{A\mid R}\arrow[d]\\
\mathcal C_{A\mid S}^{\vee}\arrow[r,"\beta^{\vee}"']&\mathcal C_{A\mid R}^{\vee}
\end{tikzcd}\tag{**}
\]
注意 $\widetilde\beta$ 未必是\Term{exact}{正合}\Term{functor}{函子}。

\textbf{3.1(b)} \Term{embedding}{嵌入}\Term{functor}{函子} $i_{A\mid R}:\widetilde{\mathcal C}_{A\mid R}\to\mathcal C_{A\mid R}^{\vee}$ 的\Term{right-derived-functor}{右导出函子}记作 $\underline H^{\bullet}_{A\mid R}$。对 $\mathcal C_{A\mid R}$ 上的每个\Term{sheaf}{层} $F$，令
\[H^{\bullet}(A\mid R,F)=[\underline H^{\bullet}_{A\mid R}(F)](A).\]
对任意 $B'\to B\in\operatorname{Cov}(B)$，定义\Term{functor}{函子}
\[H^0(\{B'\to B\},-):\mathcal C_{A\mid R}^{\vee}\longrightarrow(A\text{-mod})\]
为
\[H^0(\{B'\to B\},F)=\operatorname{Ker}\bigl(F(B')\rightrightarrows F(B'\mathbin{\times_B}B')\bigr).\]
$H^0(\{B'\to B\},-)$ 的\Term{right-derived-functor}{右导出函子}记作 $H^{\bullet}(\{B'\to B\},-)$。已知 $H^{\bullet}(\{B'\to B\},F)$ 是（Čech）\Term{semisimplicial-module}{半单纯模}
\SourcePage{82}
\[
F(B')\rightrightarrows F(B'\mathbin{\times_B}B')\mathrel{\substack{\longrightarrow\\[-.8ex]\longrightarrow\\[-.8ex]\longrightarrow}} F(B'\mathbin{\times_B}B'\mathbin{\times_B}B')\longrightarrow\cdots
\]
的\Term{cohomology-group}{上同调群}。\Term{functor}{函子} $\check{\underline H}^0_{A\mid R}:\mathcal C_{A\mid R}^{\vee}\to\mathcal C_{A\mid R}^{\vee}$ 的\Term{right-derived-functor}{右导出函子}记作 $\check{\underline H}^{\bullet}_{A\mid R}$，其中
\[
[\check{\underline H}^0_{A\mid R}(F)](B)=\varinjlim_{\{B'\to B\}\in\operatorname{Cov}(B)}H^0(\{B'\to B\},F).
\]
对 $\mathcal C_{A\mid R}$ 上的每个\Term{presheaf}{预层} $F$，令
\[
\check H^{\bullet}(A\mid R,F)=[\check{\underline H}^{\bullet}_{A\mid R}(F)](A)
=\varinjlim_{\{A'\to A\}\in\operatorname{Cov}(A)}H^{\bullet}(\{A'\to A\},F).
\]
对 $\mathcal C_{A\mid R}$ 上任意\Term{sheaf}{层} $F$，有通常的\Term{spectral-sequence}{谱序列}
\[\check H^p(A\mid R,\underline H^q_{A\mid R}(F))\Longrightarrow H^{p+q}(A\mid R,F),\]
并且注意对所有 $q>0$，有 $\check H^0(A\mid R,\underline H^q_{A\mid R}(F))=0$。特别地，
\[
H^1(A\mid R,F)=\check H^1(A\mid R,F),
\]
且序列
\[
0\to\check H^2(A\mid R,F)\to H^2(A\mid R,F)\to
\check H^1(A\mid R,\underline H^1_{A\mid R}(F))\to\check H^3(A\mid R,F)
\]
\Term{exact}{正合}。

\textbf{3.1(c)} 若 $B$ 是 $\mathcal C_{A\mid R}$ 的\Term{object}{对象}，则 2.1(a) 中的\Term{commutative-diagram}{交换图} $(*)$，连同 $\widetilde j$ 与 $j^{\vee}$ 都\Term{exact}{正合}这一事实，给出\Term{canonical-isomorphism}{典范同构}
\[H^{\bullet}(B\mid R,\widetilde jF)=[\underline H^{\bullet}_{A\mid R}(F)](B).\]
另一方面，\Term{ring-homomorphism}{环同态} $S\to R$ 诱导 2.1(a) 中的\Term{commutative-diagram}{交换图} $(**)$。这里 $\widetilde\beta$ 未必是\Term{exact}{正合}\Term{functor}{函子}。然而，$\beta$ 保持\Term{fiber-product}{纤维积}，所以得到\Term{spectral-sequence}{谱序列}
\[
\underline H^p_{A\mid R}(R^q\widetilde\beta F)\Longrightarrow\beta^{\vee}\underline H^{p+q}_{A\mid S}(F),
\quad\text{即}\quad
H^p(A\mid R,R^q\widetilde\beta F)\Longrightarrow H^{p+q}(A\mid S,F).
\]
由于有\Term{canonical-isomorphism}{典范同构}
\[\widetilde\beta=G_{A\mid R}\circ\beta^{\vee}\circ i_{A\mid S},\]
且 $\beta^{\#}=G_{A\mid R}\circ\beta^{\vee}$ 是\Term{exact}{正合}\Term{functor}{函子}，
\SourcePage{83}
故有 $R^{\bullet}\widetilde\beta F=\beta^{\#}\underline H^{\bullet}_{A\mid S}(F)$。于是上述\Term{spectral-sequence}{谱序列}可写作
\[H^p(A\mid R,\beta^{\#}\underline H^q_{A\mid S}(F))\Longrightarrow H^{p+q}(A\mid S,F).\]

\begin{SGAresult}{\Term{lemma}{引理}}{3.2}
给定 $B\in\operatorname{ob}(\mathcal C_{A\mid R})$，在 $\mathcal C_{B\mid R}$ 中选一个\Term{surjectivity}{满射}\Term{arrow}{箭头} $P\to B$，其中 $P$ 是 $R$ 上的\Term{polynomial-algebra}{多项式代数}。则对任意 $F\in\operatorname{ob}(\mathcal C_{A\mid R}^{\vee})$，自然映射
\[H^{\bullet}(\{P\to B\},F)\longrightarrow\check{\underline H}^{\bullet}(F)(B)\]
是\Term{isomorphism}{同构}。
\end{SGAresult}
\begin{proof}
只需说明自然映射
\[
H^0(\{P\to B\},F)\longrightarrow\check{\underline H}^0(F)(B)
=\varinjlim_{\{B'\to B\}\in\operatorname{Cov}(B)}H^0(\{B'\to B\},F)
\]
是\Term{isomorphism}{同构}。这显然成立，因为对任意 $\{B'\to B\}\in\operatorname{Cov}(B)$，都有一个\Term{arrow}{箭头} $\{P\to B\}\to\{B'\to B\}$。
\end{proof}

\begin{SGAresult}{推论}{3.3}
(a) 设 $A$ \Term{noetherian}{诺特}。若 $F\in\operatorname{ob}(\widetilde{\mathcal C}_{A\mid R})$ 是\Term{finite-type}{有限型}的，则对所有 $n$，$\underline H^n(F)$ 都是\Term{finite-type}{有限型}的。

(b) 对 $R$ 上任意\Term{polynomial-algebra}{多项式代数} $P$、所有 $n>0$ 及所有 $F\in\operatorname{ob}(\widetilde{\mathcal C}_{P\mid R})$，有 $H^n(P\mid R,F)=0$。

(c) $\widetilde{\mathcal C}_{A\mid R}$ 中的图表 $0\to F'\to F\to F''\to0$ \Term{exact}{正合}，当且仅当对每个作为 $R$ 上\Term{polynomial-algebra}{多项式代数}的 $P\in\operatorname{ob}(\mathcal C_{A\mid R})$，序列 $0\to F'(P)\to F(P)\to F''(P)\to0$ \Term{exact}{正合}。
\end{SGAresult}
\begin{proof}
(a) 由\Term{spectral-sequence}{谱序列}
\[\check H^p(B\mid R,\underline H^q(F))\Longrightarrow H^{p+q}(B\mid R,F),\]
以及所有 $q>0$ 时 $\check{\underline H}^0(\underline H^q(F))=0$ 这一事实，只须证明：对所有 $n$ 以及所有\Term{finite-type}{有限型}\Term{presheaf}{预层} $G$，$\check{\underline H}^n(G)$ 都是\Term{finite-type}{有限型}的。然而，$\check H^{\bullet}(B\mid R,G)=H^{\bullet}(\{P\to B\},G)$，其中 $P$ 是 $R$ 上的\Term{polynomial-algebra}{多项式代数}，且 $\{P\to B\}\in\operatorname{Cov}(B)$；而 $H^{\bullet}(\{P\to B\},G)$ 就是由\Term{complex}{复形}
\SourcePage{84}
\[\cdots G(P\mathbin{\times_B}P\mathbin{\times_B}P)\mathrel{\substack{\longrightarrow\\[-.8ex]\longrightarrow\\[-.8ex]\longrightarrow}} G(P\mathbin{\times_B}P)\rightrightarrows G(P)\]
得到的\Term{cohomology}{上同调}。若 $B$ 是\Term{finite-type}{有限型} $R$-\Term{algebra}{代数}，可取 $P$ 为\Term{finite-type}{有限型} $R$-\Term{algebra}{代数}，于是 $P\mathbin{\times_B}\cdots\mathbin{\times_B}P$ 都是\Term{finite-type}{有限型} $R$-\Term{algebra}{代数}。若 $G$ 是\Term{finite-type}{有限型}的，则所有 $G(P\mathbin{\times_B}\cdots\mathbin{\times_B}P)$ 都是\Term{finite-type}{有限型} $A$-模；由于 $A$ \Term{noetherian}{诺特}，$H^{\bullet}(\{P\to B\},G)$ 也是\Term{finite-type}{有限型} $A$-模。

(b) 设 $P$ 是 $R$ 上的\Term{polynomial-algebra}{多项式代数}。任意\Term{surjectivity}{满射} $R$-\Term{algebra}{代数}映射 $B\to P$ 都有\Term{section}{截面}，因而对所有 $i>0$，有 $H^i(\{B\to P\},G)=0$；所以对所有 $i>0$ 及 $\mathcal C_{P\mid R}$ 上所有\Term{presheaf}{预层} $G$，都有 $\check H^i(P\mid R,G)=0$。于是\Term{spectral-sequence}{谱序列}
\[\check H^p(P\mid R,\underline H^q_{P\mid R}(F))\Longrightarrow H^{p+q}(P\mid R,F)\]
\Term{degeneration}{退化}，故对所有 $n$，有 $\check H^0(P\mid R,\underline H^n(F))\simeq H^n(P\mid R,F)$。但当 $F$ 为\Term{sheaf}{层}时，对所有 $n>0$，有 $\check H^0(P\mid R,\underline H^n(F))=0$；故对所有 $n>0$ 及所有 $F\in\operatorname{ob}(\widetilde{\mathcal C}_{P\mid R})$，有 $H^n(P\mid R,F)=0$。

(c) 若 $0\to F'\to F\to F''\to0$ 是 $\widetilde{\mathcal C}_{A\mid R}$ 中的\Term{exact-sequence}{正合列}，则
\[0\to F'(B)\to F(B)\to F''(B)\to H^1(B\mid R,F')\to H^1(B\mid R,F)\to\cdots\]
\Term{exact}{正合}。若 $P\in\operatorname{ob}(\mathcal C_{A\mid R})$ 是 $R$ 上的\Term{polynomial-algebra}{多项式代数}，则由 3.3(b)，$H^1(P\mid R,F')=0$，因而
\[0\to F'\to F\to F''\to0\]
\Term{exact}{正合}。反之，令 $\mathcal P_{A\mid R}$ 为 $\mathcal C_{A\mid R}$ 的\Term{full-subcategory}{全子范畴}，其\Term{object}{对象}为
\[\operatorname{ob}(\mathcal P_{A\mid R})=\{P\in\operatorname{ob}(\mathcal C_{A\mid R})\mid P\text{ 是 }R\text{ 上的多项式代数}\},\]
赋予诱导\Term{topology}{拓扑}，并令 $p:\mathcal P_{A\mid R}\to\mathcal C_{A\mid R}$ 为包含\Term{functor}{函子}。由于 $\mathcal C_{A\mid R}$ 的每个\Term{object}{对象}都能由 $\mathcal P_{A\mid R}$ 的一个\Term{object}{对象}\Term{covering}{覆盖}，由\Term{comparison-lemma}{比较引理}（\SGAsourcegap{PDF 第84页，\Term{comparison-lemma}{比较引理}后的引注\Term{lie-bracket}{括号}原文为空，未载参考编号。}）可知 $\widetilde p:\widetilde{\mathcal C}_{A\mid R}\to\widetilde{\mathcal P}_{A\mid R}$ 是\Term{equivalence-of-categories}{范畴等价}。所以，若对每个 $P\in\operatorname{ob}(\mathcal P_{A\mid R})$，$0\to F'(P)\to F(P)\to F''(P)\to0$ 都\Term{exact}{正合}，则 $0\to\widetilde pF'\to\widetilde pF\to\widetilde pF''\to0$ 在 $\widetilde{\mathcal P}_{A\mid R}$ 中\Term{exact}{正合}，从而 $0\to F'\to F\to F''\to0$ 在 $\widetilde{\mathcal C}_{A\mid R}$ 中\Term{exact}{正合}。
\end{proof}

\SourcePage{85}
\textbf{注记 3.4。} $\mathcal P_{A\mid R}$ 中的每个\Term{surjectivity}{满射}\Term{arrow}{箭头}都有\Term{section}{截面}。立即可知，$\mathcal P_{A\mid R}$ 上的每个\Term{presheaf}{预层}实际上都是\Term{sheaf}{层}，即 $\mathcal P_{A\mid R}^{\vee}=\widetilde{\mathcal P}_{A\mid R}$。因此，\Term{restriction}{限制}\Term{functor}{函子} $\widetilde{\mathcal C}_{A\mid R}\to\mathcal P_{A\mid R}^{\vee}$ 是\Term{equivalence-of-categories}{范畴等价}。

\begin{SGAresult}{\Term{lemma}{引理}}{3.5}
考虑 $R$-\Term{algebra}{代数}图表
\[
\begin{tikzcd}&C\arrow[d]\\ B'\arrow[r,two heads,"\text{满射}"']&B\end{tikzcd}
\]
并设 $S$ 是 $R$-\Term{algebra}{代数}。

(i) 若 $\operatorname{Tor}_1^R(B,S)=0$，则\Term{canonical-map}{典范映射}
\[(B'\mathbin{\times_B}C)\otimes_R S\longrightarrow(B'\otimes_R S)\mathbin{\times_{B\otimes_R S}}(C\otimes_R S)\]
是\Term{isomorphism}{同构}。

(ii) 设 $B'$ 在 $R$ 上\Term{flat}{平坦}。若所有 $0<i<n$ 时 $\operatorname{Tor}_i^R(C,S)=0$，且所有 $0<i\leq n$ 时 $\operatorname{Tor}_i^R(B,S)=0$，则所有 $0<i<n$ 时 $\operatorname{Tor}_i^R(B'\mathbin{\times_B}C,S)=0$。
\end{SGAresult}
\begin{proof}
设 $0\to J\to B'\to B\to0$ \Term{exact}{正合}。

(i) 若 $\operatorname{Tor}_1^R(B,S)=0$，则
\[0\to J\otimes_R S\to B'\otimes_R S\to B\otimes_R S\to0\]
\Term{exact}{正合}，从而
\[0\to J\otimes_R S\to (B'\otimes_R S)\mathbin{\times_{B\otimes_R S}}(C\otimes_R S)\to C\otimes_R S\to0\]
\Term{exact}{正合}。于是得到各行\Term{exact}{正合}的\Term{commutative-diagram}{交换图}
\[
\begin{tikzcd}[column sep=small]
&J\otimes_R S\arrow[r]\arrow[d,equal]&(B'\mathbin{\times_B}C)\otimes_R S\arrow[r]\arrow[d]&C\otimes_R S\arrow[r]\arrow[d,equal]&0\\
0\arrow[r]&J\otimes_R S\arrow[r]&(B'\otimes_R S)\mathbin{\times_{B\otimes_R S}}(C\otimes_R S)\arrow[r]&C\otimes_R S\arrow[r]&0
\end{tikzcd}
\]
\SourcePage{86}
因此 $(B'\mathbin{\times_B}C)\otimes_R S\to(B'\otimes_R S)\mathbin{\times_{B\otimes_R S}}(C\otimes_R S)$ 是\Term{isomorphism}{同构}。

(ii) 由假设，$B'$ 在 $R$ 上\Term{flat}{平坦}，且对所有 $0<i\leq n$，$\operatorname{Tor}_i^R(B,S)=0$，故对所有 $0<i<n$，$\operatorname{Tor}_i^R(J,S)=0$。再由\Term{exact-sequence}{正合列} $0\to J\to B'\mathbin{\times_B}C\to C\to0$，以及所有 $0<i<n$ 时 $\operatorname{Tor}_i^R(C,S)=0$ 的假设，得到所有 $0<i<n$ 时 $\operatorname{Tor}_i^R(B'\mathbin{\times_B}C,S)=0$。
\end{proof}

\begin{SGAresult}{推论}{3.6}
(a) 设 $A,S$ 是 $R$-\Term{algebra}{代数}。若 $\operatorname{Tor}_1^R(A,S)=0$，则对任意 $G\in\operatorname{ob}(\mathcal C_{A\otimes_R S\mid S}^{\vee})$，\Term{canonical-map}{典范映射}
\[\check H^{\bullet}(A\otimes_R S\mid S,G)\longrightarrow\check H^{\bullet}(A\mid R,S^{\vee}G)\]
是\Term{isomorphism}{同构}，这里\Term{functor}{函子} $S:\mathcal C_{A\mid R}\to\mathcal C_{A\otimes_R S\mid S}$ 由 $X\mapsto X\otimes_R S$ 给出。特别地，若 $S$ 在 $R$ 上\Term{flat}{平坦}，则对每个 $G\in\operatorname{ob}(\mathcal C_{A\otimes_R S\mid S}^{\vee})$，\Term{canonical-map}{典范映射}
\[S^{\vee}\check{\underline H}^{\bullet}(G)\longrightarrow\check{\underline H}^{\bullet}(S^{\vee}G)\]
是\Term{isomorphism}{同构}。

(b) 设 $P\to B$ 是\Term{surjectivity}{满射} $R$-\Term{algebra}{代数}映射，其中 $P$ 是 $R$ 上的\Term{polynomial-algebra}{多项式代数}。若对所有 $0<i\leq n$，$\operatorname{Tor}_i^R(B,S)=0$，则对所有 $0<i<n$，
\[\operatorname{Tor}_i^R(P\mathbin{\times_B}P\mathbin{\times_B}\cdots\mathbin{\times_B}P,S)=0.\]
\end{SGAresult}
\begin{proof}
(a) 给定 $B\in\operatorname{ob}(\mathcal C_{A\mid R})$，选取 $\{P\to B\}\in\operatorname{Cov}(B)$，其中 $P$ 是 $R$ 上的\Term{polynomial-algebra}{多项式代数}。若 $\operatorname{Tor}_1^R(B,S)=0$，则由 3.5(i)，对所有 $n$，自然映射
\[(\underbrace{P\mathbin{\times_B}\cdots\mathbin{\times_B}P}_{n})\otimes_R S\longrightarrow\underbrace{(P\otimes_R S)\mathbin{\times_{B\otimes_R S}}\cdots\mathbin{\times_{B\otimes_R S}}(P\otimes_R S)}_{n}\]
是\Term{isomorphism}{同构}，因此
\[
\begin{aligned}
\check H^{\bullet}(B\mid R,S^{\vee}G)
&=H^{\bullet}(\{P\to B\},S^{\vee}G)\\
&=H^{\bullet}(\{P\otimes_R S\to B\otimes_R S\},G)
=\check H^{\bullet}(B\otimes_R S\mid S,G).
\end{aligned}
\]
若 $S$ 在 $R$ 上\Term{flat}{平坦}，则对每个 $X\in\operatorname{ob}(\mathcal C_{A\mid R})$，$\operatorname{Tor}_1^R(X,S)=0$，故
\[\check H^{\bullet}(X\otimes_R S\mid S,G)\longrightarrow\check H^{\bullet}(X\mid R,S^{\vee}G)\]
对每个这样的 $X$ 都是\Term{isomorphism}{同构}；即 $S^{\vee}\check{\underline H}^{\bullet}(G)\to\check{\underline H}^{\bullet}(S^{\vee}G)$ 是\Term{isomorphism}{同构}。

(b) 令 $P^{(m)}=\underbrace{P\mathbin{\times_B}\cdots\mathbin{\times_B}P}_{m+1}$。若 $m=0$，由于 $P$ 在 $R$ 上\Term{flat}{平坦}，无须证明。设对所有 $0<i<n$，$\operatorname{Tor}_i^R(P^{(m-1)},S)=0$。由于 $P\to B$ \Term{surjectivity}{满射}，且 $P^{(m)}=P\mathbin{\times_B}P^{(m-1)}$，由 3.5(ii)，对所有 $0<i<n$，有 $\operatorname{Tor}_i^R(P^{(m)},S)=0$。
\end{proof}

\SourcePage{87}
\begin{SGAresult}{命题}{3.7}
(a) 设 $A,S$ 是 $R$-\Term{algebra}{代数}。若对所有 $0<i\leq n$，$\operatorname{Tor}_i^R(A,S)=0$，则对 $\mathcal C_{A\otimes_R S\mid S}$ 上任意\Term{sheaf}{层} $F$，\Term{canonical-map}{典范映射}
\[H^i(A\otimes_R S\mid S,F)\longrightarrow H^i(A\mid R,\widetilde SF)\]
对所有 $i\leq n$ 都是\Term{isomorphism}{同构}，这里 $S:\mathcal C_{A\mid R}\to\mathcal C_{A\otimes_R S\mid S}$ 由 $X\mapsto X\otimes_R S$ 给出。

(b) 设 $P$ 是 $R$ 上的\Term{polynomial-algebra}{多项式代数}。则对 $\mathcal C_{P\otimes_R A\mid R}$ 上任意\Term{additive-sheaf}{加性层} $F$，\Term{canonical-map}{典范映射}
\[H^n(P\otimes_R A\mid R,F)\longrightarrow H^n(A\mid R,F)\]
对所有 $n>0$ 都是\Term{isomorphism}{同构}。

(c) 设 $S\to R\to A$ 为\Term{ring-homomorphism}{环同态}。则对 $\mathcal C_{A\mid S}$ 上任意\Term{additive-sheaf}{加性层} $F$，有\Term{exact-sequence}{正合列}
\[
\begin{aligned}
0&\to H^0(A\mid R,F_R)\to H^0(A\mid S,F)\to H^0(R\mid S,F)\\
&\to H^1(A\mid R,F_R)\to H^1(A\mid S,F)\to H^1(R\mid S,F)\to\cdots\\
&\to H^n(A\mid R,F_R)\to H^n(A\mid S,F)\to H^n(R\mid S,F)\\
&\to H^{n+1}(A\mid R,F_R)\to\cdots,
\end{aligned}
\]
这里 $F_R$ 是 $\mathcal C_{A\mid R}$ 上的\Term{sheaf}{层}，通过遗忘\Term{functor}{函子} $\mathcal C_{A\mid R}\to\mathcal C_{A\mid S}$，对所有 $X\in\operatorname{ob}(\mathcal C_{A\mid R})$ 定义为
\[F_R(X)=\operatorname{Ker}(F(X)\to F(R)).\]
\end{SGAresult}
\begin{proof}
(a) 考虑\Term{continuous-functor}{连续函子}
\[S:\mathcal C_{A\mid R}\longrightarrow\mathcal C_{A\otimes_R S\mid S},\qquad S(X)=X\otimes_R S.\]
由于 $S$ 将 $R$ 上的\Term{polynomial-algebra}{多项式代数}送到 $S$ 上的\Term{polynomial-algebra}{多项式代数}，由 3.3(c)，$\widetilde S:\widetilde{\mathcal C}_{A\otimes_R S\mid S}\to\widetilde{\mathcal C}_{A\mid R}$ 是\Term{exact}{正合}\Term{functor}{函子}，从而得到\Term{canonical-map}{典范映射}
\[S^{\vee}\underline H^{\bullet}_{A\otimes_R S\mid S}(F)\longrightarrow\underline H^{\bullet}_{A\mid R}(\widetilde SF),\]
进而得到如下\Term{spectral-sequence}{谱序列}的\Term{commutative-diagram}{交换图}：
\[
\begin{tikzcd}[column sep=small]
E_S^{p,q}(S(X))=\check H^p(X\otimes_R S\mid S,\underline H^q_{A\otimes_R S\mid S}(F))\arrow[d]\arrow[r,Rightarrow]&H^{p+q}(X\otimes_R S\mid S,F)\arrow[dd]\\
\check H^p(X\mid R,S^{\vee}\underline H^q_{A\otimes_R S\mid S}(F))\arrow[d]&\\
E_R^{p,q}(X)=\check H^p(X\mid R,\underline H^q_{A\mid R}(\widetilde SF))\arrow[r,Rightarrow]&H^{p+q}(X\mid R,\widetilde SF).
\end{tikzcd}
\]
\SourcePage{88}
注意，由 3.6(a)，只要 $\operatorname{Tor}_1^R(X,S)=0$，映射
\[E_S^{p,q}(S(X))\longrightarrow\check H^p(X\mid R,S^{\vee}\underline H^q_{A\otimes_R S\mid S}(F))\]
就是\Term{isomorphism}{同构}。现假设对所有 $q<n$，以及所有满足 $\operatorname{Tor}_i^R(X,S)=0$（原文范围写作 $0\quad i\quad n$；\SGAsourcegap{PDF 第88页，3.7(a) 证明归纳假设中指标范围的两个比较符均缺失，现存原图仅有“0 i n”。}）的 $X\in\operatorname{ob}(\mathcal C_{A\mid R})$，映射
\[[S^{\vee}\underline H^q_{A\otimes_R S\mid S}(F)](X)\longrightarrow[\underline H^q_{A\mid R}(\widetilde SF)](X)\]
是\Term{isomorphism}{同构}。由 3.6(b) 立即可知，只要 $\operatorname{Tor}_i^R(X,S)=0$ 对所有 $0<i\leq n$ 成立，$E_S^{p,q}(S(X))\to E_R^{p,q}(X)$ 对所有 $q<n$ 及所有 $p$ 都是\Term{isomorphism}{同构}。又因为 $E_S^{0,n}(S(X))\to E_R^{0,n}(X)$ 是\Term{isomorphism}{同构}（两边均为零），故只要 $\operatorname{Tor}_i^R(X,S)=0$ 对所有 $0<i\leq n$ 成立，$H^n(X\otimes_R S\mid S,F)\to H^n(X\mid R,\widetilde SF)$ 就是\Term{isomorphism}{同构}；即在同样条件下，
\[[S^{\vee}\underline H^i_{A\otimes_R S\mid S}(F)](X)\longrightarrow[\underline H^i_{A\mid R}(\widetilde SF)](X)\]
对所有 $i\leq n$ 都是\Term{isomorphism}{同构}。

(b) 设 $P$ 是 $R$ 上的\Term{polynomial-algebra}{多项式代数}，令 $f:\mathcal C_{A\mid R}\to\mathcal C_{A\otimes_R P\mid R}$ 为由 $f(X)=X\otimes_R P$ 给出的\Term{continuous-functor}{连续函子}。设 $P'\to X$ 是 $\mathcal C_{A\mid R}$ 中的\Term{surjectivity}{满射}\Term{arrow}{箭头}，其中 $P'$ 是 $R$ 上的\Term{polynomial-algebra}{多项式代数}。于是 $P'\otimes_R P\to X\otimes_R P$ 是 $\mathcal C_{A\otimes_R P\mid R}$ 中的\Term{surjectivity}{满射}\Term{arrow}{箭头}，且 $P'\otimes_R P$ 仍是 $R$ 上的\Term{polynomial-algebra}{多项式代数}。由于 $P$ 在 $R$ 上\Term{flat}{平坦}，由 3.5(i)，对每个 $G\in\operatorname{ob}(\mathcal C_{A\otimes_R P\mid R}^{\vee})$，有
\[
\begin{aligned}
\check H^{\bullet}(X\mid R,f^{\vee}G)
&=H^{\bullet}(\{P'\to X\},f^{\vee}G)\\
&=H^{\bullet}(\{P'\otimes_R P\to X\otimes_R P\},G)
=\check H^{\bullet}(X\otimes_R P\mid R,G).
\end{aligned}
\]
于是 $\widetilde f:\widetilde{\mathcal C}_{A\otimes_R P\mid R}\to\widetilde{\mathcal C}_{A\mid R}$ 将\Term{acyclic}{无上同调}\Term{object}{对象}送到\Term{acyclic}{无上同调}\Term{object}{对象}。另一方面，$f$ 将 $R$ 上的\Term{polynomial-algebra}{多项式代数}送到 $R$ 上的\Term{polynomial-algebra}{多项式代数}，所以由 3.2(c)，$\widetilde f:\widetilde{\mathcal C}_{A\otimes_R P\mid R}\to\widetilde{\mathcal C}_{A\mid R}$ 是\Term{exact}{正合}\Term{functor}{函子}。从而对每个 $F\in\operatorname{ob}(\widetilde{\mathcal C}_{A\otimes_R P\mid R})$，$H^{\bullet}(A\otimes_R P\mid R,F)\to H^{\bullet}(A\mid R,\widetilde fF)$ 是\Term{isomorphism}{同构}。现在，对每个 $X\in\operatorname{ob}(\mathcal C_{A\mid R})$，有\Term{canonical-map}{典范映射}
\[(\widetilde fF)(X)=F(X\otimes_R P)\longrightarrow F(X)\oplus F(P),\]
因而有\Term{canonical-map}{典范映射} $\widetilde fF\to F\oplus F(P)$，这里 $F(P)$ 表示\Term{constant-sheaf}{常值层}（严格地说，右端的 $F$ 指 $F$ 在 $\mathcal C_{A\mid R}$ 上的\Term{restriction}{限制}）。若 $F$ 是\Term{additive-sheaf}{加性层}，则对每个 $X\in\operatorname{ob}(\mathcal C_{A\mid R})$，
\[(\widetilde fF)(X)=F(X\otimes_R P)\xrightarrow{\sim}F(X)\oplus F(P)\]
是\Term{isomorphism}{同构}，因而 $\widetilde fF\xrightarrow{\sim}F\oplus F(P)$ 是\Term{isomorphism}{同构}。于是
\SourcePage{89}
\[H^n(A\mid R,\widetilde fF)\xrightarrow{\sim}H^n(A\mid R,F\oplus F(P))=H^n(A\mid R,F)\]
对所有 $n>0$ 成立，因为 $F(P)$ 是\Term{constant-sheaf}{常值层}，所以是\Term{flasque-sheaf}{松层}。故只要 $F$ 是 $\mathcal C_{A\otimes_R P\mid R}$ 上的\Term{additive-sheaf}{加性层}，$H^n(A\otimes_R P\mid R,F)\to H^n(A\mid R,F)$ 对所有 $n>0$ 都是\Term{isomorphism}{同构}。

(c)（由 Rinehart 告知。）考虑遗忘\Term{functor}{函子} $\beta:\mathcal C_{A\mid R}\to\mathcal C_{A\mid S}$。有\Term{spectral-sequence}{谱序列}
\[E^{p,q}=H^p(A\mid R,\beta^{\#}\underline H^q_{A\mid S}(F))\Longrightarrow H^{p+q}(A\mid S,F),\]
其中 $\beta^{\#}=G_{A\mid R}\circ\beta^{\vee}$，而 $G_{A\mid R}:\mathcal C_{A\mid R}^{\vee}\to\widetilde{\mathcal C}_{A\mid R}$ 是“\Term{associated-sheaf}{伴随层}”\Term{functor}{函子}。设 $F$ 是\Term{additive-sheaf}{加性层}。我们断言，对所有 $q>0$，$\beta^{\#}\underline H^q_{A\mid S}(F)$ 是\Term{constant-sheaf}{常值层}。事实上，若 $P$ 是 $R$ 上的\Term{polynomial-algebra}{多项式代数}，则对某个 $S$ 上的\Term{polynomial-algebra}{多项式代数} $P'$，有 $P=R\otimes_S P'$；从而由 3.7(b)，
\[[\beta^{\#}\underline H^q_{A\mid S}(F)](P)=H^q(P\mid S,F)=H^q(R\otimes_S P'\mid S,F)\longrightarrow H^q(R\mid S,F)\]
对所有 $q>0$ 都是\Term{isomorphism}{同构}，而当 $q=0$ 时\Term{surjectivity}{满射}。由 3.3(c)，\Term{canonical-map}{典范映射} $\beta^{\#}\underline H^q_{A\mid S}(F)\to H^q(R\mid S,F)$（右边是\Term{constant-sheaf}{常值层}）对所有 $q>0$ 都是\Term{isomorphism}{同构}，当 $q=0$ 时\Term{surjectivity}{满射}。因此，对所有 $p,q>0$，有 $E^{p,q}=H^p(A\mid R,\beta^{\#}\underline H^q_{A\mid S}(F))=0$，上述\Term{spectral-sequence}{谱序列}遂给出\Term{exact-sequence}{正合列}
\[0\to E^{1,0}\to H^1\to E^{0,1}\to E^{2,0}\to H^2\to E^{0,2}\to E^{3,0}\to H^3\to\cdots.\tag{*}\]
这里注意，对所有 $q>0$，有 $E^{0,q}=[\beta^{\#}\underline H^q_{A\mid S}(F)](A)\xrightarrow{\sim}H^q(R\mid S,F)$，且 $E^{p,0}=H^p(A\mid R,\widetilde\beta F)$。另一方面，\Term{exact-sequence}{正合列} $0\to F_R\to\widetilde\beta F\to F(R)\to0$ 给出\Term{exact-sequence}{正合列}
\[
\begin{aligned}
0\to H^0(A\mid R,F_R)&\to\underbrace{H^0(A\mid R,\widetilde\beta F)}_{=H^0(A\mid S,F)}
\to\underbrace{H^0(A\mid R,F(R))}_{=H^0(R\mid S,F)}\\
&\to H^1(A\mid R,F_R)\to H^1(A\mid R,\widetilde\beta F)\to0.
\end{aligned}\tag{**}
\]
结合\Term{exact-sequence}{正合列} $(*)$ 与 $(**)$，即得所求\Term{exact-sequence}{正合列}。
\end{proof}

\SourcePage{90}
\begin{SGAresult}{推论}{3.8}
(a) 若 $A,B$ 是 $R$-\Term{algebra}{代数}，且对所有 $0<i\leq n$，$\operatorname{Tor}_i^R(A,B)=0$，则对任意\Term{additive-sheaf}{加性层} $F\in\operatorname{ob}(\widetilde{\mathcal C}_{A\otimes_R B\mid R})$，\Term{canonical-map}{典范映射}
\[H^i(A\otimes_R B\mid R,F)\longrightarrow H^i(A\mid R,F)\oplus H^i(B\mid R,F)\]
对所有 $i\leq n$ 都是\Term{isomorphism}{同构}。

(b) 设 $A$ 是 $R$-\Term{algebra}{代数}，且对所有 $i>0$，$\operatorname{Tor}_i^R(A,A)=0$。若 $A\otimes_R A\to A$ 是\Term{isomorphism}{同构}，则对任意\Term{additive-sheaf}{加性层} $F\in\operatorname{ob}(\widetilde{\mathcal C}_{A\mid R})$ 以及所有 $i$，有 $H^i(A\mid R,F)=0$。
\end{SGAresult}
\begin{proof}
(a) 设 $F\in\operatorname{ob}(\widetilde{\mathcal C}_{A\otimes_R B\mid R})$ 是\Term{additive}{加性}的。由 3.7(c)，\Term{ring-homomorphism}{环同态} $R\to A\to A\otimes_R B$ 给出\Term{exact-sequence}{正合列}
\[\cdots\to H^i(A\otimes_R B\mid A,F_A)\to H^i(A\otimes_R B\mid R,F)\to H^i(A\mid R,F)\to\cdots.\tag{*}\]
现考虑\Term{commutative-diagram}{交换图}
\[
\begin{tikzcd}
H^i(A\otimes_R B\mid A,F_A)\arrow[r]\arrow[d]&H^i(A\otimes_R B\mid R,F)\arrow[d]\\
H^i(B\mid R,\widetilde AF_A)\arrow[r,"\sim"']&H^i(B\mid R,F),
\end{tikzcd}
\]
其中\Term{functor}{函子} $A:\mathcal C_{B\mid R}\to\mathcal C_{A\otimes_R B\mid A}$ 由 $X\mapsto A\otimes_R X$ 给出，映射 $H^i(B\mid R,\widetilde AF_A)\to H^i(B\mid R,F)$ 由映射 $\widetilde AF_A\to F$ 诱导；由于 $F$ 是\Term{additive}{加性}的，后一个映射是\Term{isomorphism}{同构}。注意，由\Term{functoriality}{函子性}，上述图的\Term{commutativity}{交换性}归结为 $i=0$ 的情形。若 $\operatorname{Tor}_i^R(A,B)=0$ 对所有 $0<i\leq n$ 成立，则由 3.7(a)，$H^i(A\otimes_R B\mid A,F_A)\to H^i(B\mid R,\widetilde AF_A)$ 对所有 $i\leq n$ 都是\Term{isomorphism}{同构}，故复合映射
\[H^i(A\otimes_R B\mid A,F_A)\to H^i(A\otimes_R B\mid R,F)\to H^i(B\mid R,F),\]
以及（由\Term{symmetry}{对称性}）
\SourcePage{91}
\[H^i(A\otimes_R B\mid B,F_B)\to H^i(A\otimes_R B\mid R,F)\to H^i(A\mid R,F)\]
对所有 $i\leq n$ 都是\Term{isomorphism}{同构}。这说明\Term{exact-sequence}{正合列} $(*)$ \Term{resolution}{分解}为\Term{split-exact-sequence}{分裂短正合列}
\[0\to H^i(A\otimes_R B\mid A,F_A)\to H^i(A\otimes_R B\mid R,F)\to H^i(A\mid R,F)\to0\]
（$i\leq n$），并且
\[H^i(A\otimes_R B\mid R,F)\longrightarrow H^i(A\mid R,F)\oplus H^i(B\mid R,F)\]
对所有 $i\leq n$ 都是\Term{isomorphism}{同构}。

(b) 因为对所有 $i>0$，$\operatorname{Tor}_i^R(A,A)=0$，由 3.8(a)，
\[H^i(A\otimes_R A\mid R,F)\longrightarrow H^i(A\mid R,F)\oplus H^i(A\mid R,F)\]
对所有 $i$ 都是\Term{isomorphism}{同构}。另一方面，由假设 $A\otimes_R A\to A$ 是\Term{isomorphism}{同构}，故 $H^i(A\mid R,F)\to H^i(A\otimes_R A\mid R,F)$ 对所有 $i$ 都是\Term{isomorphism}{同构}。这就是说，\Term{diagonal-map}{对角映射} $H^i(A\mid R,F)\to H^i(A\mid R,F)\oplus H^i(A\mid R,F)$ 对所有 $i$ 都是\Term{isomorphism}{同构}，即对所有 $i$，$H^i(A\mid R,F)=0$。
\end{proof}

\begin{SGAresult}{推论}{3.9}
设 $S$ 是 $R$-\Term{algebra}{代数} $A$ 的\Term{multiplication}{乘法}\Term{subset}{子集}。若 $F\in\operatorname{ob}(\widetilde{\mathcal C}_{S^{-1}A\mid R})$ 是\Term{additive}{加性}的，且 $F_A$ 也是\Term{additive}{加性}的，则\Term{canonical-map}{典范映射}
\[H^i(S^{-1}A\mid R,F)\longrightarrow H^i(A\mid R,F)\]
对所有 $i$ 都是\Term{isomorphism}{同构}；这里 $F_A\in\operatorname{ob}(\widetilde{\mathcal C}_{S^{-1}A\mid A})$ 由 $F_A(X)=\operatorname{Ker}(F(X)\to F(A))$ 定义，其中 $X\in\operatorname{ob}(\mathcal C_{S^{-1}A\mid A})$。
\end{SGAresult}
\begin{proof}
由于 $F$ 是\Term{additive}{加性}的，\Term{ring-homomorphism}{环同态} $R\to A\to S^{-1}A$ 给出\Term{exact-sequence}{正合列}
\[\cdots\to H^n(S^{-1}A\mid A,F_A)\to H^n(S^{-1}A\mid R,F)\to H^n(A\mid R,F)\to\cdots.\]
另一方面，$S^{-1}A\otimes_A S^{-1}A\to S^{-1}A$ 是\Term{isomorphism}{同构}，且由假设 $F_A$ 是\Term{additive}{加性}的，故由 3.8(b)，对所有 $n$，$H^n(S^{-1}A\mid A,F_A)=0$；因而 $H^n(S^{-1}A\mid R,F)\to H^n(A\mid R,F)$ 对所有 $n$ 都是\Term{isomorphism}{同构}。
\end{proof}

\SourcePage{92}
现设 $A$ 是 $R$-\Term{algebra}{代数}。$A$-模 $M$ 通过 $X\mapsto\operatorname{Der}_R(X,M)$ 定义\Term{additive-sheaf}{加性层} $\operatorname{Der}_R(-,M)\in\operatorname{ob}(\widetilde{\mathcal C}_{A\mid R})$。现在作如下定义。
\begin{SGAresult}{定义}{3.10}
设 $A$ 是 $R$-\Term{algebra}{代数}。对每个 $A$-模 $M$，令
\[D^i(A\mid R,M)=H^i(A\mid R,\operatorname{Der}_R(-,M)).\]
\end{SGAresult}
我们要显式计算低维情形，即 $i=1,2$ 时的 $D^i(A\mid R,M)$。在 $\mathcal C_{A\mid R}$ 中选取\Term{surjectivity}{满射}\Term{arrow}{箭头} $P\to A$，其中 $P$ 是 $R$ 上的\Term{polynomial-algebra}{多项式代数}，并设 $0\to I\to P\to A\to0$ \Term{exact}{正合}。则有：
\begin{SGAresult}{\Term{lemma}{引理}}{3.10}
\[
\begin{aligned}
D^1(A\mid R,M)&=\operatorname{Coker}(\operatorname{Der}_R(P,M)\to\operatorname{Hom}_P(I,M)),\\
D^2(A\mid R,M)&=H^1(A\mid R,\underline H^1_{A\mid R}(\operatorname{Der}_R(-,M))).
\end{aligned}
\]
\end{SGAresult}
\begin{proof}
考虑\Term{simplicial-algebra}{单纯代数}
\[\cdots\mathrel{\substack{\longrightarrow\\[-.8ex]\longrightarrow\\[-.8ex]\longrightarrow}}P_A^{(2)}\mathrel{\substack{\longrightarrow\\[-.8ex]\longrightarrow\\[-.8ex]\longrightarrow}}P_A^{(1)}\rightrightarrows P_A^{(0)},\tag{*}\]
其中 $P_A^{(n)}=\underbrace{P\mathbin{\times_A}\cdots\mathbin{\times_A}P}_{n+1}$。对每个 $n$，有\Term{exact-sequence}{正合列}
\[0\to I^{(n)}\to P_A^{(n)}\to A\to0,\]
其中 $I^{(n)}=\underbrace{I\times\cdots\times I}_{n+1}$，并且容易看出，对所有 $n$，序列
\[0\to\operatorname{Hom}_A(\overline I^{(n)}/\Delta(\overline I),M)\to\operatorname{Der}_R(P_A^{(n)},M)\xrightarrow{\operatorname{Der}(\Delta,M)}\operatorname{Der}_R(P,M)\to0\tag{a}\]
\Term{exact}{正合}，这里 $\overline I=I/I^2$，$\Delta$ 表示\Term{diagonal-map}{对角映射}。还注意到\Term{exact-sequence}{正合列}
\[0\to\operatorname{Hom}_A(\overline I^{(n)}/\Delta(\overline I),M)\to\operatorname{Hom}_A(\overline I^{(n)},M)\xrightarrow{\operatorname{Hom}(\Delta,M)}\operatorname{Hom}_A(\overline I,M)\to0.\tag{b}\]
\Term{exact-sequence}{正合列} (a)、(b) 分别给出相应的\Term{simplicial-module}{单纯} $A$-模\Term{exact-sequence}{正合列}。由于 $\{\overline I\to0\}$、$\{\overline I\to\overline I\}$ 都有\Term{section}{截面}，故
\[H^i(\{\overline I\to0\},\operatorname{Hom}_A(-,M))=H^i(\{\overline I\to\overline I\},\operatorname{Hom}_A(-,M))=0\]
\SourcePage{93}
对所有 $i>0$ 成立，随即得出：
\[0\to\operatorname{Der}_R(A,M)\to\operatorname{Der}_R(P,M)\to\operatorname{Hom}_A(\overline I,M)\to H^1(\{P\to A\},\operatorname{Der}_R(-,M))\to0\]
\Term{exact}{正合}，且对所有 $i>1$，
\[H^i(\{P\to A\},\operatorname{Der}_R(-,M))=0.\]
由 3.1(b)，
\[
\begin{aligned}
D^1(A\mid R,M)&=\check H^1(A\mid R,\operatorname{Der}_R(-,M))
=H^1(\{P\to A\},\operatorname{Der}_R(-,M))\\
&=\operatorname{Coker}(\operatorname{Der}_R(P,M)\to\operatorname{Hom}_A(\overline I,M))\\
&=\operatorname{Coker}(\operatorname{Der}_R(P,M)\to\operatorname{Hom}_P(I,M)),\\
D^2(A\mid R,M)&=H^1(A\mid R,\underline H^1_{A\mid R}(\operatorname{Der}_R(-,M))).
\end{aligned}
\]
\end{proof}
现选取\Term{free-module}{自由} $P$-模 $F$ 及 $P$-\Term{linear}{线性}映射 $\gamma:F\to P$，使 $\gamma(F)=I$。与 $P$-\Term{linear}{线性}映射 $\gamma:F\to P$ \Term{adjunction}{伴随}有一个 \Term{koszul-complex}{Koszul 复形}，滥用记号，将它简记为 $K_I$。

\begin{SGAresult}{\Term{lemma}{引理}}{3.11}
\[D^2(A\mid R,M)=\operatorname{Coker}(H^1(K_I,M)\to\operatorname{Hom}_A(H_1(K_I),M)).\]
\end{SGAresult}
\begin{proof}
有
\[
\begin{aligned}
D^2(A\mid R,M)&=H^1(A\mid R,\underline H^1_{A\mid R}(\operatorname{Der}_R(-,M)))\\
&=\operatorname{Ker}\bigl(H^1(P_A^{(1)}\mid R,M)\mathrel{\substack{\longrightarrow\\[-.7ex]\longrightarrow\\[-.7ex]\longrightarrow}}H^1(P_A^{(2)}\mid R,M)\bigr),
\end{aligned}
\]
因为 $[\underline H^1_{A\mid R}(\operatorname{Der}_R(-,M))](P_A^{(0)})=[\underline H^1_{A\mid R}(\operatorname{Der}_R(-,M))](P)=0$。
设 $0\to K\to F\to I\to0$ \Term{exact}{正合}。若以 $P[F]$ 表示由 $F$ 生成的 $P$ 上的\Term{symmetric-algebra}{对称代数}，则 $P[F]$ 和 $P[F]\otimes_P P[F]$ 都是 $R$ 上的\Term{polynomial-algebra}{多项式代数}，并得到\Term{surjectivity}{满射}\Term{arrow}{箭头}
\[
\begin{aligned}
P[F]&\longrightarrow P_A^{(1)}=P\mathbin{\times_A}P,&X&\longmapsto(0,\gamma(X)),\\
P[F]\otimes_P P[F]&\longrightarrow P_A^{(2)}=P\mathbin{\times_A}P\mathbin{\times_A}P,
&X\otimes1&\longmapsto(0,\gamma(X),0),\quad1\otimes X\longmapsto(0,0,\gamma(X)),
\end{aligned}
\]
其中 $X\in F$。令 $J=\operatorname{Ker}(P[F]\to P_A^{(1)})$ 和 $J'=\operatorname{Ker}(P[F]\otimes_P P[F]\to P_A^{(2)})$，则
\SourcePage{94}
\[
\begin{aligned}
J&=(F)\cap(F_\gamma)=(K)+(F\cdot F_\gamma),\\
J'&=(F\otimes1+1\otimes F)\cap(F_\gamma\otimes1+1\otimes F)\cap(F\otimes1+1\otimes F_\gamma)\\
&=(K\otimes1+1\otimes K)+(F\cdot F_\gamma\otimes1+1\otimes F\cdot F_\gamma)+(F\otimes F),
\end{aligned}
\]
其中 $F_\gamma$ 是由 $\{X-\gamma(X)\mid X\in F\}$ 生成的\Term{ideal}{理想}。映射 $\pi_i:P_A^{(2)}\to P_A^{(1)}$（$i=0,1,2$）\Term{liftable}{可提升}为 $\widetilde\pi_i:P[F]\otimes_P P[F]\to P[F]$，其中对每个 $X\in F$，
\[
\begin{array}{ll}
\widetilde\pi_0:\;X\otimes1\mapsto\gamma(X)-X,&1\otimes X\mapsto X,\\
\widetilde\pi_1:\;X\otimes1\mapsto0,&1\otimes X\mapsto X,\\
\widetilde\pi_2:\;X\otimes1\mapsto X,&1\otimes X\mapsto0.
\end{array}
\]
直接计算可见，$\widetilde\pi_0-\widetilde\pi_1+\widetilde\pi_2$ 将 $K\otimes1+1\otimes K+1\otimes F\cdot F_\gamma$ 送到零，而将 $(F\cdot F_\gamma\otimes1+F\otimes F)$ 映满到 $(F\cdot F_\gamma)$。现对任意 $D\in\operatorname{Der}_R(P[F]\otimes_P P[F],M)$，由于 $I\cdot M=0$，有 $D(F\cdot F_\gamma\otimes1+F\otimes F)=0$。因此
\[
\begin{aligned}
D^2(A\mid R,M)
&=\operatorname{Ker}\bigl(D^1(P_A^{(1)}\mid R,M)\mathrel{\substack{\longrightarrow\\[-.7ex]\longrightarrow\\[-.7ex]\longrightarrow}}D^1(P_A^{(2)}\mid R,M)\bigr)\\
&=\{[f]\in D^1(P_A^{(1)}\mid R,M)\mid f(F\cdot F_\gamma)=0\}\\
&=\operatorname{Coker}\bigl(\operatorname{Der}_R(P[F],M)\to\operatorname{Hom}_{P[F]}((F)\cap(F_\gamma)/(F\cdot F_\gamma),M)\bigr)\\
&=\operatorname{Coker}\bigl(\operatorname{Der}_P(P[F],M)\to\operatorname{Hom}_{P[F]}((F)\cap(F_\gamma)/(F)\cdot(F_\gamma),M)\bigr).
\end{aligned}
\]
现在，序列
\[
\begin{aligned}
0\to(F)\to P[F]\xrightarrow{P[0]}P\to0,\\
0\to(F_\gamma)\to P[F]\xrightarrow{P[\gamma]}P\to0
\end{aligned}
\]
\Term{exact}{正合}，因此
\[(F)\cap(F_\gamma)/(F)\cdot(F_\gamma)=\operatorname{Tor}^{P[F]}_1(P,P_\gamma)=H_1(K_I),\]
其中 $K_I$ 是与 $P$-\Term{linear}{线性}映射 $F\xrightarrow{\gamma}I\subset P$ \Term{adjunction}{伴随}的 \Term{koszul-complex}{Koszul 复形}。所以有
\[
\begin{aligned}
D^2(A\mid R,M)&=\operatorname{Coker}(\operatorname{Hom}_P(F,M)\to\operatorname{Hom}_P(H_1(K_I),M))\\
&=\operatorname{Coker}(H^1(K_I,M)\to\operatorname{Hom}_A(H_1(K_I),M)).
\end{aligned}
\]
\end{proof}
现在，我们可以将已经得到的各项结果汇总起来，陈述所需的全部基本事实。

\SourcePage{95}
\begin{SGAresult}{\Term{theorem}{定理}}{3.12}
(1) 若 $0\to M'\to M\to M''\to0$ 是 $A$-模的\Term{exact-sequence}{正合列}，则序列
\[
\begin{aligned}
0&\to D^0(A\mid R,M')\to D^0(A\mid R,M)\to D^0(A\mid R,M'')\to D^1(A\mid R,M')\\
&\to D^1(A\mid R,M)\to D^1(A\mid R,M'')\to\cdots\to D^n(A\mid R,M')\\
&\to D^n(A\mid R,M)\to D^n(A\mid R,M'')\to\cdots
\end{aligned}
\]
\Term{exact}{正合}。

(2) 设 $S\to R\to A$ 为\Term{ring-homomorphism}{环同态}。则对任意 $A$-模 $M$，有\Term{exact-sequence}{正合列}
\[
\begin{aligned}
0&\to D^0(A\mid R,M)\to D^0(A\mid S,M)\to D^0(R\mid S,M)\to D^1(A\mid R,M)\\
&\to D^1(A\mid S,M)\to D^1(R\mid S,M)\to\cdots\to D^n(A\mid R,M)\\
&\to D^n(A\mid S,M)\to D^n(R\mid S,M)\to\cdots.
\end{aligned}
\]

(3) 设 $A,S$ 是 $R$-\Term{algebra}{代数}，且对所有 $i>0$，$\operatorname{Tor}_i^R(A,S)=0$。则对任意 $A\otimes_R S$-模 $M$，\Term{canonical-map}{典范映射}
\[D^i(A\otimes_R S\mid S,M)\longrightarrow D^i(A\mid R,M)\]
对所有 $i$ 都是\Term{isomorphism}{同构}。特别地，若 $A$ 在 $R$ 上\Term{flat}{平坦}，则对所有 $i$ 及所有 $R$-\Term{algebra}{代数} $S$，有
\[D^i(A\otimes_R S\mid S,M)\xrightarrow{\sim}D^i(A\mid R,M).\]

(4) 设 $S$ 是 $A$ 的\Term{multiplication}{乘法}\Term{subset}{子集}。则对任意 $S^{-1}A$-模 $N$，\Term{canonical-map}{典范映射}
\[D^i(S^{-1}A\mid R,N)\longrightarrow D^i(A\mid R,N)\]
对所有 $i$ 都是\Term{isomorphism}{同构}。因此，对任意 $A$-模 $M$ 及所有 $i$，有\Term{canonical-isomorphism}{典范同构}
\[D^i(S^{-1}A\mid R,S^{-1}M)\xrightarrow{\sim}S^{-1}D^i(A\mid R,M).\]

(5) 若 $A,B$ 是 $R$-\Term{algebra}{代数}，且对所有 $0<i\leq n$，$\operatorname{Tor}_i^R(A,B)=0$，则对任意 $A\otimes_R B$-模 $M$，\Term{canonical-map}{典范映射}
\[D^i(A\otimes_R B\mid R,M)\longrightarrow D^i(A\mid R,M)\oplus D^i(B\mid R,M)\]
对所有 $i\leq n$ 都是\Term{isomorphism}{同构}。特别地，若对所有 $i>0$，$\operatorname{Tor}_i^R(A,A)=0$，且
\SourcePage{96}
\[\Delta:\operatorname{Spec}(A)\longrightarrow\operatorname{Spec}(A)\mathbin{\times_{\operatorname{Spec}(R)}}\operatorname{Spec}(A)\]
是\Term{open-immersion}{开浸入}（例如 $A/R$ \Term{etale}{平展}），则对所有 $i>0$ 及所有 $A$-模 $M$，有 $D^i(A\mid R,M)=0$。

(6) 设 $R$ \Term{noetherian}{诺特}，$A$ 是 $R$ 上\Term{essentially-finite-type}{本质有限型}的 $R$-\Term{algebra}{代数}，$E$ 是\Term{flat}{平坦} $A$-模。则对所有 $i\geq0$ 及所有 $A$-模 $M$，有\Term{canonical-isomorphism}{典范同构}
\[E\otimes_A D^i(A\mid R,M)\xrightarrow{\sim}D^i(A\mid R,E\otimes_A M).\]

(7) 若 $R$ \Term{noetherian}{诺特}且 $A$ 是\Term{essentially-finite-type}{本质有限型} $R$-\Term{algebra}{代数}，则只要 $M$ 是\Term{finite-type}{有限型} $A$-模，$D^i(A\mid R,M)$ 对所有 $i$ 都是\Term{finite-type}{有限型} $A$-模。

(8) 设 $P$ 是 $R$ 上的\Term{polynomial-algebra}{多项式代数}，并设 $0\to I\to P\to A\to0$ \Term{exact}{正合}。则
\[
\begin{aligned}
D^0(A\mid R,M)&=\operatorname{Der}_R(A,M),\\
D^1(A\mid R,M)&=\operatorname{Coker}(\operatorname{Der}_R(P,M)\to\operatorname{Hom}_A(I/I^2,M))\\
&=\{\text{以 }M\text{ 为核的 }A\text{ 的 }R\text{-代数扩张的同构类}\},\\
D^2(A\mid R,M)&=\operatorname{Coker}(H^1(K_I,M)\to\operatorname{Hom}_A(H_1(K_I),M)).
\end{aligned}
\]
\end{SGAresult}
\begin{proof}
(1) 若 $0\to M'\to M\to M''\to0$ 是 $A$-模的\Term{exact-sequence}{正合列}，则对 $R$ 上任意\Term{polynomial-algebra}{多项式代数} $P$，$0\to\operatorname{Der}_R(P,M')\to\operatorname{Der}_R(P,M)\to\operatorname{Der}_R(P,M'')\to0$ \Term{exact}{正合}，故由 3.3(c)，$0\to\operatorname{Der}_R(-,M')\to\operatorname{Der}_R(-,M)\to\operatorname{Der}_R(-,M'')\to0$ 是\Term{sheaf}{层}的\Term{exact-sequence}{正合列}。

(2) $\operatorname{Der}_S(-,M)$ 是 $\mathcal C_{A\mid S}$ 上的\Term{additive-sheaf}{加性层}，且对每个 $B\in\operatorname{ob}(\mathcal C_{A\mid R})$，有
\[\operatorname{Der}_R(B,M)=\operatorname{Ker}(\operatorname{Der}_S(B,M)\to\operatorname{Der}_S(R,M)).\]
故结论由 3.7(c) 得到。

(3) 令 $S:\mathcal C_{A\mid R}\to\mathcal C_{A\otimes_R S\mid S}$ 为由 $X\mapsto X\otimes_R S$ 给出的\Term{functor}{函子}。于是 $\widetilde S\operatorname{Der}_S(-,M)=\operatorname{Der}_R(-,M)$，故结论由 3.7(a) 得到。

\SourcePage{97}
(4) 由 3.9 得到。事实上，$F=\operatorname{Der}_R(-,N)$ 是 $\mathcal C_{S^{-1}A\mid R}$ 上的\Term{additive-sheaf}{加性层}，而 $F_A\in\operatorname{ob}(\widetilde{\mathcal C}_{S^{-1}A\mid A})$ 也是\Term{additive}{加性}的，因为
\[F_A(X)=\operatorname{Ker}(F(X)\to F(A))=\operatorname{Ker}(\operatorname{Der}_R(X,N)\to\operatorname{Der}_R(A,N))=\operatorname{Der}_A(X,N),\]
即 $F_A=\operatorname{Der}_A(-,N)$。

(5) 第一个断言由 3.8(a) 得到。现设对所有 $i>0$，$\operatorname{Tor}_i^R(A,A)=0$，且
\[\Delta:\operatorname{Spec}(A)\longrightarrow\operatorname{Spec}(A)\mathbin{\times_{\operatorname{Spec}(R)}}\operatorname{Spec}(A)\]
是\Term{open-immersion}{开浸入}；即对 $A$ 的所有\Term{maximal-ideal}{极大理想} $\mathfrak m$，$(A\otimes_R A)_{\varphi^{-1}(\mathfrak m)}\to A_{\mathfrak m}$ 是\Term{isomorphism}{同构}，这里 $\varphi:A\otimes_R A\to A$ 为\Term{multiplication}{乘法}映射。于是第一个断言与 3.12(4) 合起来说明，对每个\Term{maximal-ideal}{极大理想} $\mathfrak m$，\Term{diagonal-map}{对角映射}
\[\Delta_{\mathfrak m}:D^i(A\mid R,M)_{\mathfrak m}\longrightarrow D^i(A\mid R,M)_{\mathfrak m}\oplus D^i(A\mid R,M)_{\mathfrak m}\]
是\Term{isomorphism}{同构}，因而
\[\Delta:D^i(A\mid R,M)\longrightarrow D^i(A\mid R,M)\oplus D^i(A\mid R,M)\]
是\Term{isomorphism}{同构}，遂得对所有 $i$，有 $D^i(A\mid R,M)=0$。

(6) 由 3.12(4)，可以设 $A$ 是\Term{finite-type}{有限型} $R$-\Term{algebra}{代数}。现考虑 $\mathcal C_{A\mid R}^{\vee}$ 中的映射
\[E\otimes_A\operatorname{Der}_R(-,M)\longrightarrow\operatorname{Der}_R(-,E\otimes_A M).\]
由于 $E$ 在 $A$ 上\Term{flat}{平坦}，只要 $B\in\operatorname{ob}(\mathcal C_{A\mid R})$ 在 $R$ 上\Term{finite-type}{有限型}，就有\Term{isomorphism}{同构}
\[E\otimes_A\operatorname{Der}_R(B,M)\xrightarrow{\sim}\operatorname{Der}_R(B,E\otimes_A M).\]
由 3.3(b)，
\[D^i(A\mid R,E\otimes_A\operatorname{Der}_R(-,M))\longrightarrow D^i(A\mid R,E\otimes_A M)\]
是\Term{isomorphism}{同构}。另一方面，\Term{functor}{函子} $E\otimes_A-:(A\text{-mod})\to(A\text{-mod})$ \Term{exact}{正合}，所以
\[E\otimes_A D^i(A\mid R,M)=D^i(A\mid R,E\otimes_A\operatorname{Der}_R(-,M)),\]
从而得到结论。

(7) 由 3.3(a) 得到。

\SourcePage{98}
(8) 由 3.10 和 3.11，只需说明
\[\operatorname{Coker}(\operatorname{Der}_R(P,M)\to\operatorname{Hom}_A(I/I^2,M))=\operatorname{Exalcomm}(A\mid R,M),\]
这里 $\operatorname{Exalcomm}(A\mid R,M)$ 表示以 $M$ 为核的 $A$ 的 $R$-\Term{algebra}{代数}\Term{group-extension}{扩张}的\Term{isomorphism-class}{同构类}所成的\Term{set}{集合}。设 $0\to M\to A'\to A\to0$ 是以 $M$ 为核的 $A$ 的一个 $R$-\Term{algebra}{代数}\Term{group-extension}{扩张}。因为 $P$ 是 $R$ 上的\Term{polynomial-algebra}{多项式代数}，得到\Term{commutative-diagram}{交换图}
\[
\begin{tikzcd}
0\arrow[r]&I\arrow[r]\arrow[d]&P\arrow[r]\arrow[d]&A\arrow[r]\arrow[d,equal]&0\\
0\arrow[r]&M\arrow[r]&A'\arrow[r]&A\arrow[r]&0,
\end{tikzcd}
\]
其中 $I\to M$ 由\Term{group-extension}{扩张} $A'$ 唯一确定到相差 $\operatorname{Der}_R(P,M)$ 中的一个元素。因此得到映射
\[\Phi:\operatorname{Exalcomm}(A\mid R,M)\longrightarrow\operatorname{Coker}(\operatorname{Der}_R(P,M)\to\operatorname{Hom}_A(I/I^2,M)).\]
反之，$P$-\Term{linear}{线性}映射 $I\to M$ 通过\Term{pushout}{推出}给出一个\Term{group-extension}{扩张}：
\[
\begin{tikzcd}
0\arrow[r]&I\arrow[r]\arrow[d]&P\arrow[r]\arrow[d]&A\arrow[r]\arrow[d,equal]&0\\
0\arrow[r]&M\arrow[r]&A'\arrow[r]&A\arrow[r]&0.
\end{tikzcd}
\]
显然，两个相差 $\operatorname{Der}_R(P,M)$ 中一个元素的 $P$-\Term{linear}{线性}映射 $I\to M$ 给出\Term{isomorphic}{同构}的\Term{group-extension}{扩张}；因此得到映射
\[\Psi:\operatorname{Coker}(\operatorname{Der}_R(P,M)\to\operatorname{Hom}_A(I/I^2,M))\longrightarrow\operatorname{Exalcomm}(A\mid R,M).\]
直接可见，$\Phi$ 和 $\Psi$ 是互逆的过程。
\end{proof}

\begin{SGAresult}{推论}{3.13}
(a) $R$-\Term{algebra}{代数} $A$ 在 $R$ 上\Term{formally-smooth}{形式光滑}，当且仅当对每个 $A$-模 $M$，$D^1(A\mid R,M)=0$。

(b) 设 $R$ \Term{noetherian}{诺特}，$A$ 是\Term{finite-type}{有限型} $R$-\Term{algebra}{代数}。则 $A$ 是 $R$ 上的\Term{relative-complete-intersection}{相对完全交}，当且仅当对每个 $A$-模 $M$，$D^2(A\mid R,M)=0$。

(c) 设 $R$ \Term{noetherian}{诺特}，$A$ 是\Term{finite-type}{有限型} $R$-\Term{algebra}{代数}。则对 $A$ 的任意\Term{multiplication}{乘法}\Term{subset}{子集} $S$，以及所有 $i\geq0$，有\Term{canonical-isomorphism}{典范同构}
\[D^i(S^{-1}A\mid R,S^{-1}M)\xrightarrow{\sim}S^{-1}D^i(A\mid R,M).\]

\SourcePage{99}
(d) 设 $R$ \Term{noetherian}{诺特}，$A$ 是\Term{essentially-finite-type}{本质有限型}的\Term{local-algebra}{局部} $R$-\Term{algebra}{代数}。以 $\widehat A$ 表示 $A$ 的 $\mathfrak m$-进\Term{completion}{完备化}，其中 $\mathfrak m$ 为 $A$ 的\Term{maximal-ideal}{极大理想}。则对任意\Term{finite-type}{有限型} $A$-模 $E$，有\Term{canonical-isomorphism}{典范同构}
\[D^1(\widehat A\mid R,\widehat A\otimes_A E)\xrightarrow{\sim}\widehat A\otimes_A D^1(A\mid R,E).\]
\end{SGAresult}
\begin{proof}
(a) 由 $D^1(A\mid R,M)=\operatorname{Exalcomm}(A\mid R,M)$ 立即得到。

(b) 设 $P$ 是 $R$ 上\Term{finite}{有限}个变量的\Term{polynomial-algebra}{多项式代数}，且 $0\to I\to P\to A\to0$ \Term{exact}{正合}。必要时\Term{localization}{局部化}后，可以假定 $I$ 由一个 $P$-\Term{regular-sequence}{正则序列}生成，从而 $H_1(K_I)=0$，因此
\[D^2(A\mid R,M)=\operatorname{Coker}(H^1(K_I,M)\to\operatorname{Hom}_A(H_1(K_I),M))=0.\]
反之，设对每个 $A$-模 $M$，$D^2(A\mid R,M)=0$。
设 $0\to I\to P\to A\to0$ \Term{exact}{正合}，其中 $P$ 是 $R$ 上\Term{finite}{有限}个变量的\Term{polynomial-algebra}{多项式代数}。有
\[0=D^2(A\mid R,M)=\operatorname{Coker}(H^1(K_I,M)\to\operatorname{Hom}_A(H_1(K_I),M)),\]
即对所有 $A$-模 $M$，$H^1(K_I,M)\to\operatorname{Hom}_A(H_1(K_I),M)$ \Term{surjectivity}{满射}。选取\Term{exact-sequence}{正合列} $0\to K\to F\to I\to0$，其中 $F$ 是\Term{finite-type}{有限型}\Term{free-module}{自由} $P$-模。注意 $H_1(K_I)=K/K_0$，其中 $K_0=\operatorname{Im}(\bigwedge^2F\to\bigwedge^1F)$。由 $D^2(A\mid R,H_1(K_I))=0$，即
\[H^1(K_I,H_1(K_I))\longrightarrow\operatorname{Hom}_A(H_1(K_I),H_1(K_I))\quad\text{满射},\]
可知 $0\to K/K_0\to F/IF\to I/I^2\to0$ \Term{exact}{正合}，并作为 $A$-模\Term{split}{分裂}，特别地，$I/I^2$ 是\Term{projective-module}{投射} $A$-模。因此，必要时将 $A$ \Term{localization}{局部化}，可以假定 $I/I^2$ 是\Term{free-module}{自由} $A$-模，且 $F/IF\to I/I^2$ 是\Term{isomorphism}{同构}。于是 $0\to K/K_0\to F/IF\to I/I^2\to0$ \Term{exact}{正合}且作为 $A$-模\Term{split}{分裂}这一事实说明 $H_1(K_I)=K/K_0=0$。这意味着 $I$ 由 $P$-\Term{regular-sequence}{正则序列}生成，即 $A$ 是 $R$ 上的\Term{relative-complete-intersection}{相对完全交}。

(c) 由 3.12(4) 和 (6) 立即得到。

(d) 首先由 3.12(6)，
\[D^i(A\mid R,\widehat A\otimes_A E)\xrightarrow{\sim}\widehat A\otimes_A D^i(A\mid R,E)\]
对所有 $i$ 都是\Term{isomorphism}{同构}，所以只须说明
\SourcePage{100}
\[D^1(\widehat A\mid R,\widehat A\otimes_A E)\longrightarrow D^1(A\mid R,\widehat A\otimes_A E)\]
是\Term{isomorphism}{同构}。我们的证明基于用\Term{algebra}{代数}\Term{group-extension}{扩张}描述 $D^1$。设
\[0\to\widehat A\otimes_A E\to A'\xrightarrow{\varphi}\widehat A\to0\]
是一个 $R$-\Term{algebra}{代数}\Term{group-extension}{扩张}。$\widehat A\otimes_A E$ 是\Term{finite-type}{有限型} $\widehat A$-模这一事实立即说明，$A'$ 是\Term{local-algebra}{局部} $R$-\Term{algebra}{代数}，且关于 $\mathfrak m'$-进\Term{topology}{拓扑}\Term{complete}{完备}，其中 $\mathfrak m'$ 是 $A'$ 的\Term{maximal-ideal}{极大理想}。因此，若有 $R$-\Term{algebra}{代数}映射 $\psi:A\to A'$ 使 $\varphi\circ\psi=\operatorname{id}_A$，则 $\psi$ 唯一延拓为 $\widehat\psi:\widehat A\to A'$，使 $\varphi\circ\widehat\psi=\operatorname{id}_{\widehat A}$。这证明了
\[D^1(\widehat A\mid R,\widehat A\otimes_A E)\longrightarrow D^1(A\mid R,\widehat A\otimes_A E)\]
是\Term{injectivity}{单射}。现设 $0\to\widehat A\otimes_A E\to B\xrightarrow{\varphi}A\to0$ 是一个 $R$-\Term{algebra}{代数}\Term{group-extension}{扩张}。首先断言：$B$ 的 $\mathfrak m_B$-进\Term{topology}{拓扑}在 $\widehat E=\widehat A\otimes_A E$ 上诱导的\Term{topology}{拓扑}，等于 $\widehat E$ 作为 $\widehat A$-模的\Term{topology}{拓扑}。事实上，$A$ 在 $R$ 上\Term{essentially-finite-type}{本质有限型}，故可在 $B$ 中找到\Term{subalgebra}{子代数} $B_0$，使 $B_0$ 是 $R$ 上\Term{essentially-finite-type}{本质有限型}的\Term{local-algebra}{局部} $R$-\Term{algebra}{代数}，并且 $\varphi(B_0)=A$。令 $E_0=\widehat E\cap B_0$，分别以 $\mathfrak m_B,\mathfrak m_0,\mathfrak m$ 表示 $B,B_0,A$ 的\Term{maximal-ideal}{极大理想}。于是 $\mathfrak m_B=\mathfrak m_0+\widehat E$，从而对所有 $n$，有 $\mathfrak m_B^{n+1}=\mathfrak m_0^{n+1}+\mathfrak m^n\widehat E$。由于 $B_0$ \Term{noetherian}{诺特}，由 Artin--Rees \Term{theorem}{定理}，存在整数 $r>0$，使对所有 $n\geq r$，
\[
\begin{aligned}
\widehat E\cap\mathfrak m_B^{n+1}
&=\mathfrak m^n\widehat E+\widehat E\cap\mathfrak m_0^{n+1}
=\mathfrak m^n\widehat E+E_0\cap\mathfrak m_0^{n+1}\\
&=\mathfrak m^n\widehat E+\mathfrak m^{n+1-r}(E_0\cap\mathfrak m_0^r)
\subset\mathfrak m^{n+1-r}\widehat E.
\end{aligned}
\]
这证明了关于 $\widehat E$ 上这两个\Term{topology}{拓扑}的断言。特别地，$\widehat E=\widehat A\otimes_A E$ 关于由 $B$ 的 $\mathfrak m_B$-进\Term{topology}{拓扑}诱导的\Term{topology}{拓扑}\Term{complete}{完备}，因而得到各行\Term{exact}{正合}的\Term{commutative-diagram}{交换图}
\[
\begin{tikzcd}
0\arrow[r]&\widehat E\arrow[r]\arrow[d,equal]&B\arrow[r]\arrow[d]&A\arrow[r]\arrow[d]&0\\
0\arrow[r]&\widehat E\arrow[r]&\widehat B\arrow[r]&\widehat A\arrow[r]&0,
\end{tikzcd}
\]
其中 $\widehat B$ 表示 $B$ 的 $\mathfrak m_B$-进\Term{completion}{完备化}。这证明 $D^1(\widehat A\mid R,\widehat E)\to D^1(A\mid R,\widehat E)$ \Term{surjectivity}{满射}，证明完毕。
\end{proof}

\SourcePage{101}
\begin{SGAresult}{推论}{3.14}
设 $R$ \Term{noetherian}{诺特}，$A$ 是\Term{finite-type}{有限型} $R$-\Term{algebra}{代数}，且除一个\Term{closed-point}{闭点} $x\in\operatorname{Spec}(A)$ 外处处\Term{smooth}{光滑}。则有\Term{isomorphism}{同构}
\[D^1(A\mid R,A)\simeq D^1(A_x\mid R,A_x)\simeq D^1(\widehat A_x\mid R,\widehat A_x).\]
\end{SGAresult}
\begin{proof}
由于 $A$ 在除 $x\in\operatorname{Spec}(A)$ 外的各点均在 $R$ 上\Term{smooth}{光滑}，$D^1(A\mid R,A)$ 是\Term{finite-type}{有限型} $A$-模，其\Term{support}{支集}包含于 $\{x\}$；因此
\[D^1(A\mid R,A)=D^1(A\mid R,A)_x=\widehat{D^1(A\mid R,A)_x}.\]
结论遂由 3.13 得到。
\end{proof}

\textbf{注记 3.15。} 设 $R$ 是\Term{topological-ring}{拓扑环}，$A$ 是\Term{topological-algebra}{拓扑} $R$-\Term{algebra}{代数}。对任意 $A$-模 $M$，定义
\[D^i_{\mathrm{top}}(A\mid R,E)=\varinjlim_{\alpha}D^i(A_\alpha\mid R_\alpha,A_\alpha\otimes_A E),\]
其中 $A_\alpha=A/J_\alpha$，$R_\alpha=R/I_\alpha$，且 $J_\alpha,I_\alpha$ 分别在 $A,R$ 中为开\Term{ideal}{理想}。于是，说 $A$ 是\Term{formally-smooth}{形式光滑}的\Term{topological-algebra}{拓扑} $R$-\Term{algebra}{代数}，意味着对每个 $A$-模 $E$，都有 $D^1_{\mathrm{top}}(A\mid R,E)=0$。

现设 $X$ 是环 $R$ 上的\Term{scheme}{概形}，$E$ 是\Term{quasi-coherent}{拟凝聚} $X$-模。我们通过 $X$ 上的\Term{presheaf}{预层}，定义 $X$ 上的\Term{sheaf}{层} $\underline D^i_{X\mid R}(E)=\underline D^i(X\mid R,E)$；对每个\Term{affine-open}{仿射开集} $U\subset X$，该\Term{presheaf}{预层}取值为
\[[\underline D^i_{X\mid R}(E)](U)=D^i(\Gamma(U,\mathcal O_X)\mid R,\Gamma(U,E)).\]
若 $R$ \Term{noetherian}{诺特}且 $X$ 在 $R$ 上\Term{finite-type}{有限型}，则性质 3.13(a) 说明，$\underline D^i_{X\mid R}(E)$ 实际上可以\Term{canonical}{典范}地构造为\Term{quasi-coherent}{拟凝聚} $X$-模。

在不致混淆时，简记 $\underline D^i_{X\mid R}=\underline D^i_{X\mid R}(\mathcal O_X)$。3.13 的所有陈述现在都可以改用这些\Term{quasi-coherent-sheaf}{拟凝聚层} $\underline D^i_{X\mid R}(E)$ 来表述，留待读者自行完成。下面列出与我们目的最相关的若干性质。

\begin{SGAresult}{推论}{3.15}
设 $R$ \Term{noetherian}{诺特}，$X$ 是\Term{finite-type}{有限型} $R$-\Term{scheme}{概形}。则有：
\SourcePage{102}
(a) 对任意\Term{coherent}{凝聚} $X$-模 $E$，所有 $i$ 时 $\underline D^i_{X\mid R}(E)$ 仍是\Term{coherent}{凝聚} $X$-模；并且对所有 $i>0$，有 $\operatorname{Supp}\underline D^i_{X\mid R}(E)\subset X^{\mathrm{Sing}}$，其中
\[X^{\mathrm{Sing}}=\{x\in X\mid X\text{ 在点 }x\in X\text{ 处于 }R\text{ 上不光滑}\}.\]

(b) 若 $X$ 在 $R$ 上\Term{flat}{平坦}，则对任意 $R$-\Term{algebra}{代数} $S$ 及所有 $i$，有\Term{canonical-isomorphism}{典范同构}
\[\underline D^i_{X\otimes_R S\mid S}(E)\xrightarrow{\sim}\underline D^i_{X\mid R}(E).\]

(c) 若 $X$ 是 $R$ 上的\Term{relative-complete-intersection}{相对完全交}，则对每个\Term{surjectivity}{满射}\Term{ring-homomorphism}{环同态} $S\to R$，有\Term{exact-sequence}{正合列}
\[0\to\underline D^1_{X\mid R}(E)\to\underline D^1_{X\mid S}(E)\to\underline{\operatorname{Hom}}_{\mathcal O_X}(I/I^2\otimes_R\mathcal O_X,E)\to0,\]
其中 $I=\operatorname{Ker}(S\to R)$。若进一步假设 $R$ 是\Term{local-ring}{局部环}，\Term{residue-field}{剩余域}为 $k$，且 $X$ 在 $R$ 上\Term{flat}{平坦}，则有\Term{exact-sequence}{正合列}
\[0\to\underline D^1_{X_0\mid k}(\mathcal O_{X_0})\to\underline D^1_{X\mid S}(\mathcal O_{X_0})\to D^1(R\mid S,k)\otimes_k\mathcal O_{X_0}\to0,\]
其中 $X_0=X\otimes_R k$。
\end{SGAresult}
\begin{proof}
(a) 由 3.12(6) 和 (8) 得到；(b) 则由 3.12(3) 得到。

(c) 由于 $X$ 是 $R$ 上的\Term{relative-complete-intersection}{相对完全交}，由 3.12(10)，对所有 $i\geq2$，有 $\underline D^i_{X\mid R}(E)=0$，故由 3.12(2)，序列
\[
\begin{aligned}
0\to\underline D^0_{X\mid R}(E)&\to\underline D^0_{X\mid S}(E)\to\underline D^0_{R\mid S}(E)\\
&\to\underline D^1_{X\mid R}(E)\to\underline D^1_{X\mid S}(E)\to\underline D^1_{R\mid S}(E)\to0
\end{aligned}
\]
\Term{exact}{正合}，其中 $\underline D^i_{R\mid S}(E)$ 是 $X$ 上的\Term{quasi-coherent-sheaf}{拟凝聚层}，对每个\Term{affine-open}{仿射开子集} $U\subset X$，由
\[[\underline D^i_{R\mid S}(E)](U)=D^i(R\mid S,E(U))\]
定义。特别地，由于 $S\to R$ \Term{surjectivity}{满射}，$\underline D^0_{R\mid S}(E)=0$；并且由 3.12(7)，
\[\underline D^1_{R\mid S}(E)=\underline{\operatorname{Hom}}_{\mathcal O_X}(I/I^2\otimes_R\mathcal O_X,E),\]
遂得第一个\Term{exact-sequence}{正合列}。若 $X$ 在 $R$ 上\Term{flat}{平坦}，则对所有 $i$ 以及所有\Term{quasi-coherent}{拟凝聚} $X_0$-模 $E$，有
\SourcePage{103}
\[\underline D^i_{X_0\mid k}(E)\xrightarrow{\sim}\underline D^i_{X\mid R}(E),\]
特别地，
\[\underline D^1_{X_0\mid k}(\mathcal O_{X_0})\xrightarrow{\sim}\underline D^1_{X\mid R}(\mathcal O_{X_0}),\]
第二个\Term{exact-sequence}{正合列}由此得到。
\end{proof}

\textbf{注记 3.16。} 我们定义的\Term{quasi-coherent-sheaf}{拟凝聚层} $\underline D^i_{X\mid R}(E)$ 应当作为某个联系\Term{local-scheme}{局部}\Term{object}{对象}与整体\Term{object}{对象}的\Term{spectral-sequence}{谱序列}中的 $E_2^{0,i}$ 出现。L.~Illusie 已经得到了这样一种理论（《Lecture Notes in Mathematics》第 239 卷）。他的方法借助“\Term{homotopical-algebra}{同伦代数}”，推广了 M.~André 与 D.~Quillen 的方法，而非借助“\Term{cohomological-algebra}{上同调代数}”。看来，这里采用的方法可以借助\Term{scheme}{概形}\Term{category}{范畴}上的 Grothendieck \Term{cohomology}{上同调}作整体化，所用的适当\Term{topology}{拓扑}推广了\Term{affine}{仿射}\Term{category}{范畴}上的\Term{topology}{拓扑}。
